% Study-2 prose. Every number is a macro generated from the artifacts (scripts/report_study2.py);
% every qualitative statement is checked by scripts/check_claims.py.

\newcommand{\StudyTwoSummary}{%
Of the \STwoNHyp\ frozen decisions, \STwoNSupported\ were supported. Apart from the
recalibration contrast C5b, every superiority contrast favoured EOST-M on at least seven
of the eight benchmarks and no competitor was significantly better on any benchmark; four
of the five superiority hypotheses that were not supported (C2, C4, C7, C8) missed the
frozen threshold by one benchmark, with four of eight significant and none against.
CRUXEval-O, with only \STwoTargetsMin\ independent clusters, did not reach significance in
any contrast.}

\newcommand{\StudyTwoReplicationText}{%
With independent targets the study-1 result therefore holds: the gain over plug-in Fisher
CAT was positive on every benchmark, larger with four calibrators (median
\SecRTwoNFourMed\%, \SecRTwoNFourFav\ benchmarks significant) and smaller with sixteen
(\SecRTwoNSixteenMed\%, none significant), and of the same size for deterministic
(\STwoRTwoDet\%) and recorded aggregate verdicts (\STwoRTwoAgg\%). Plug-in Fisher CAT's 90\%
intervals covered the full-bank score only \STwoCovmfipi\ of the time, against
\STwoCoveigmg\ for EOST-M.}

\newcommand{\StudyTwoRaschText}{%
Integration therefore beat constraint on held-out prediction, as the post hoc probe of
study~1 suggested, but by a small margin, of similar size for deterministic and aggregate
verdicts (\STwoCOneDet\% and \STwoCOneAgg\%); the margin was larger with sixteen calibrators
than with four (median \SecCOneNSixteenMed\% against \SecCOneNFourMed\%), where a single
discrimination becomes a visible misspecification. On the score, the two tests were close
(median \SecEostVsRaschIAEMed\% in favour of EOST-M, no benchmark significant), and the
Rasch-type intervals were again better calibrated before recalibration (coverage
\STwoCovmfiOnepr\ against \STwoCoveigmg; Section~\ref{sec:s2cov}).}

\newcommand{\StudyTwoReductionText}{%
The adaptive gain over Random:mg grew with the panel (median \SecCFourNFourMed\% with four
calibrators, \SecCFourNSixteenMed\% with sixteen) and was similar for deterministic and
aggregate verdicts (\STwoCFourDet\% and \STwoCFourAgg\%). The ATLAS-style test had the
largest score error of all model-based policies and covered the score only
\STwoCovatlas\ of the time: with eight calibrators its item parameters are plug-in estimates,
as in plug-in Fisher CAT, and its randomesque choice among the five most informative items
trades efficiency for exposure control (ATLAS itself calibrates on thousands of models; our
implementation follows its model and selection rule, not its code).
tinyBenchmarks and MINCE choose their items without looking at the target. With eight
calibrators, tinyBenchmarks' K-Means cannot always find as many anchors as its budget: its
item embeddings come from the panel's plug-in fit and can only take as many distinct values
as there are distinct response patterns among the items, which is far fewer than the $2^8$
possible ones when many items are solved by all or by none of the calibrators. At a budget of
100 items, fewer anchors were available in some cells of \TinyShortHundredBanks\ of the 8
benchmarks and in \TinyShortHundredMax\% of the cells of the most affected one
(Table~\ref{tab:tinyaudit}). A review of the code during the run found this, and a
post hoc audit decided before the results were seen shows that it does not change C8: at the
50-item budget it affected at most \TinyShortFiftyMax\% of the cells of a benchmark, and
restricting C8 to full-budget cells, or comparing with EOST-M after the number of items
tinyBenchmarks actually used, again gave \TinyFullFav\ significant benchmarks under the
frozen Holm family, as in the frozen analysis (median \TinyFullMed\%). Because the frozen code did
not record how many anchors were used, they were counted post hoc with the same seeds; the
recomputed errors equal the frozen ones exactly (maximum difference \TinyReproDiff). MINCE's own stopping rule
asked for a median of \STwoMinceNstarMed\ items and then had a score error of
\STwoMinceErrNstar\ points, more than EOST-M after 100 items (\STwoErreigmgHundred).
Table~\ref{tab:reduction} also lists tinyBenchmarks' p-IRT and anchor-only estimators; C8
concerns gp-IRT, the estimator its authors recommend.}

\newcommand{\StudyTwoCoverageText}{%
Model-based intervals remained too narrow with eight calibrators
(Figure~\ref{fig:study2cov}): EOST-M covered the full-bank score \STwoCoveigmg\ of the time,
the PDM Bayesian CAT \FbBmiCov\ (after the burn-in correction), the Rasch-type test \STwoCovmfiOnepr\ and plug-in
Fisher CAT \STwoCovmfipi, whereas Random:mg reached \STwoCovrandommg\ and Naive's Wilson
interval \STwoCovnaive. Leave-one-calibrator-out recalibration raised EOST-M's coverage to
\STwoCoveigmgRR\ (from \STwoCFiveaMin\ on \STwoEostRCovMinBank\ to \STwoCFiveaMax), inside
$[0.85,0.95]$ on \STwoCFiveaIn\ benchmarks, one short of the frozen criterion (C5a
\STwoCFiveaVerdict). It did not improve the interval score, which penalizes width as well as
misses (C5b \STwoCFivebVerdict; significantly worse on HellaSwag and WinoGrande): the wider
intervals cost about as much as the misses they removed. Recalibrated EOST-M intervals had a
much better interval score than Naive's Wilson intervals (C6 \STwoCSixVerdict, median
\STwoCSixMed\%), because at a coverage of \STwoCoveigmgRR\ they were much narrower (mean width
\STwoWidtheigmgRR\ against \STwoWidthnaive). The same recalibration brought the Rasch-type
test inside $[0.85,0.95]$ on all eight benchmarks (\STwoCovmfiOneprRR\ on average, between
\STwoRaschRCovMin\ and \STwoRaschRCovMax), but with intervals about twice as wide as
recalibrated EOST-M's (\STwoWidthmfiOneprRR) and a worse interval score on
\SecEostRVsRaschRISFav\ benchmarks (none better). Recalibration needs calibrators to serve as
pseudo-targets and helped more as the panel grew: EOST-M's recalibrated coverage was
\STwoCovEostRNFour, \STwoCovEostRNEight\ and \STwoCovEostRNSixteen\ at $N=4$, 8 and 16, and
its interval score was worse than the raw interval's at $N=4$ (\STwoISEostRNFour\ against
\STwoISEostNFour) but better at $N=16$ (\STwoISEostRNSixteen\ against \STwoISEostNSixteen).
Mixing one uniformly drawn item in four into the adaptive test raised coverage only to
\STwoCovmixmg.}

\newcommand{\StudyTwoRobustText}{%
\emph{Temporal validation.} On the \TempTargets\ SWE-bench Verified submissions of the latest
quarter, calibrated only on earlier submissions, \TempNSig\ of the \TempNTests\ frozen
contrasts favoured EOST-M at a Holm-adjusted $p<0.05$ within the component and none favoured
a competitor; only the recalibration contrast C5b was not significant
($p=\TempCFivebP$). EOST-M's held-out log loss was \TempRTwoMed\% lower than plug-in Fisher
CAT's (\TempRTwoWins\ of \TempRTwoN\ targets) and \TempCOneMed\% lower than the Rasch-type
test's, and its score error \TempCFourMed\% lower than Random:mg's. Coverage was
\TempCovEost\ raw and \TempCovEostR\ recalibrated (Rasch-type: \TempCovRasch\ and
\TempCovRaschR; Naive: \TempCovNaive).
\par
\emph{Synthetic banks.} Where the fitted 2PL is correct, recalibration moved coverage away
from the nominal level (\SynTwoplCovEost\ raw, \SynTwoplCovEostR\ recalibrated at $N=8$), as
it did on the bank with guessing (\SynGuessCovEost\ and \SynGuessCovEostR); it moved coverage
towards it only on the bank with a second ability dimension (\SynTwodimCovEost\ and
\SynTwodimCovEostR). Recalibration from a handful of pseudo-targets is noisy, and it pays only
when the model's own intervals are clearly too narrow.
\par
\emph{Statistical conventions.} As in study~1, the direction of a significant contrast is
read from the mean difference while the test is rank-based; across both studies \SignsDisc\
of the \SignsSig\ Holm-significant contrasts had a mean difference and a rank-biserial
correlation of opposite signs.}

\newcommand{\StudyTwoIntroMsgTwo}{%
In study~2, with one target per family cluster, EOST-M also predicted held-out verdicts
better than the regularized Rasch-type test (median \STwoCOneMed\%, \STwoCOneFav\ of 8
benchmarks significant), and Bayesian CATs that also integrate the calibrators' abilities
were not significantly different from it.}

\newcommand{\StudyTwoFBText}{%
A review of the code during the run, before any result was inspected, found that the
frozen burn-in of 10 sweeps was too short on some benchmarks: the translation move is
restricted to shifts that keep every sampled difficulty on the grid, so on HellaSwag, where
many difficulties sit near the edge of the grid, it rarely acted and the mean calibrator
ability drifted by a median of \FbDriftFrozenMax\ during the retained sweeps (median over
benchmarks \FbDriftFrozenMedian). Both Bayesian-CAT arms were re-run post hoc with a burn-in
of 60 sweeps on the same targets, panels, holdouts and chains, which brought the drift on
HellaSwag down to \FbDriftCorrMax, close to the other benchmarks (median
\FbDriftCorrMedian). The correction did not change the conclusion (Table~\ref{tab:fbfix});
Section~\ref{sec:s2follow} examines the sampler's convergence with long chains and better moves.
With either burn-in, the PDM Bayesian CAT was not significantly different from EOST-M on
any benchmark, in held-out log loss or in score error; its log loss was higher by a median
of \FbBmiMed\% with Bayesian maximum information and \FbEigMed\% with EIG (frozen:
\FbFrozenBmiMed\% and \FbFrozenEigMed\%), within $\pm1\%$ on only \FbBmiEq\ (BMI) and
\FbEigEq\ (EIG) benchmarks, and its intervals covered \FbBmiCov\ of the time against \FbEostCov\ for EOST-M.
The PDM Bayesian CAT follows \citet{vanderlinden2020fast} and averages each
item's response function over the posterior draws separately, which treats the shared
uncertainty about the location of the calibration as independent item-level noise
(Lemma~\ref{lem:exact}). The exploratory joint fully Bayesian CAT, which keeps this
uncertainty shared by mixing over draws of the whole calibration, was the closest to EOST-M
with EIG selection
(median \FbjEigMed\%, within $\pm1\%$ on \FbjEigEq\ benchmarks; with Bayesian maximum
information \FbjBmiMed\%), as we had expected before running it. With eight calibrators,
integrating their abilities and the hyperprior on top of the item parameters therefore made
no significant difference; equivalence within $\pm1\%$ was not shown, and on a few
benchmarks the bootstrap intervals still allow differences of about one per cent.}
\newcommand{\StudyTwoFBDetails}{%
The post hoc re-analysis of the frozen arm reproduces the equivalence test of
Table~\ref{tab:equiv} exactly, and with the frozen burn-in the correction runs reproduce the
frozen rows exactly (maximum difference 0 on the first target
of every benchmark), so the corrected arms differ from the frozen ones only by the burn-in.
At $N=4$ (Table~\ref{tab:fbfixfour}) the corrected PDM Bayesian CAT was within $\pm1\%$ of
EOST-M on \FbBmiEqNFour\ (BMI) and \FbEigEqNFour\ (EIG) benchmarks, with median differences of
\FbBmiMedNFour\% and \FbEigMedNFour\%, \FbBmiFbSigNFour\ and \FbEigFbSigNFour\ benchmarks
significantly in its favour and \FbBmiEostSigNFour\ and \FbEigEostSigNFour\ in EOST-M's, and
coverage \FbBmiCovNFour\ (EOST-M: \STwoCovEostNFour); the location of the calibration is more
uncertain with four calibrators, and the chains' drift was correspondingly larger. The sampler is not expensive: with eight calibrators the frozen chain
(42 sweeps) took a median of \STwoSecFb~s per panel, against \STwoSecCal~s for EOST-M's whole
calibration and \STwoSecLoco~s for the leave-one-calibrator-out recalibration.}
\newcommand{\DiscussionWhyNotFB}{%
A fully Bayesian adaptive test is the principled answer to calibration uncertainty, and
\citet{fink2025accounting} found it the most accurate correction in simulations with
hundreds of calibration responses per item. With eight calibrators it did not help
measurably: integrating the calibrators' abilities and the hyperprior changed held-out
prediction, score error and coverage by no significant amount, although equivalence within
$\pm1\%$ was not shown either. The form that averages each item over the posterior draws
separately (the PDM Bayesian CAT) was, if anything, slightly worse than EOST-M in study~2,
which would fit its treating the shared uncertainty about the location of the calibration as
independent item noise; but with the better samplers of Section~\ref{sec:s2follow} neither form
was consistently worse, so we do not read this difference as a property of the posterior-draw
approximation.
By Lemma~\ref{lem:exact}, EOST-M is the fully Bayesian model with the calibrators' abilities
and the hyperprior held fixed; in this regime what matters is the item-level uncertainty that
both integrate rather than the panel-level uncertainty that only the fully Bayesian model
adds, as far as our data can tell. A
fully Bayesian test also needs a Markov chain per panel, whose tuning is a source of error in
its own right: our frozen burn-in was too short on some benchmarks, and with eight calibrators
the posterior has a ridge in the scale of the calibration along which a standard Gibbs sampler
barely moves, so that even long chains did not converge for the scale hyperparameters; an exact
sampler with collapsed moves came close but did not meet our convergence criterion either
(Section~\ref{sec:s2follow}). EOST-M is the simpler default. Two leads remain open. The
better samplers left held-out log loss essentially unchanged but suggest a possible advantage
of Bayesian integration for estimating the score: the near-converged exact sampler was ahead of
EOST-M in score error on most of the 48 cells it was run on, but on only about half of the
benchmarks, too few to establish it. And panel-level uncertainty may matter more when the
calibrators' abilities are poorly determined, for instance with short banks; we did not test
this.}

\newcommand{\DiscussionRemedies}{%
The post hoc probes of study~1 suggested that constraining the discriminations to a common
value is almost as good as integrating their uncertainty. Study~2 tested this under a frozen
plan with independent targets. Integration was better on held-out prediction (C1
\STwoCOneVerdict, median \STwoCOneMed\%), more so with sixteen calibrators, where one common
discrimination becomes a visible misspecification; the score errors of the two tests were
close; and the constrained model's raw intervals were better calibrated because they are
wider, not because its estimates are more accurate (Section~\ref{sec:strong}). The
one-parameter model integrated over its item posterior sat in between (C2, median
\STwoCTwoMed\%). The practical reading is that a small panel must not be trusted with
item-specific discriminations: integrating their uncertainty is preferable when prediction
and score accuracy matter most, and constraining them gives a simpler and more conservative
test.}

\newcommand{\DiscussionUncertainty}{%
The review of study~1 asked whether the under-coverage of model-based intervals could be
fixed. Study~2 tried four remedies under a frozen plan: integrating the calibrators'
abilities and the hyperprior (the Bayesian CATs), leave-one-calibrator-out
recalibration, mixing uniformly drawn items into the adaptive test, and recalibrating a
constrained model. None gave nominal coverage with intervals as narrow as EOST-M's. The
Bayesian CATs' intervals and the mixed test stayed below 0.80; recalibration brought
EOST-M to \STwoCoveigmgRR\ on average, inside the frozen band on \STwoCFiveaIn\ of 8
benchmarks, with an interval score no better than before; the recalibrated Rasch-type test
reached \STwoCovmfiOneprRR\ with intervals twice as wide. Random items scored by the same
model reached \STwoCovrandommg\ without any recalibration. The exploratory coverage
decomposition of study~1 suggests that the shortfall arises when adaptive selection meets a
misspecified model, and the prior predictive checks locate a misspecification in the
dependence structure of real systems' verdicts (Section~\ref{sec:eb}). Recalibration is
therefore not yet a validated remedy: it raised coverage, but neither frozen criterion for it
was met, it made the interval score worse on two benchmarks, and on the well-specified
synthetic bank it moved coverage away from the nominal level. Practitioners should report
model-based intervals together with the panel size and their expected under-coverage, may
add recalibrated intervals as a sensitivity analysis, and should use a Wilson interval from
uniformly sampled items when nominal coverage matters more than width.}

\newcommand{\DiscussionReduction}{%
If the aim is to evaluate the next system with a fraction of a benchmark and only a handful
of calibration systems exist, which method should be used? In study~2, EOST-M's score error
after 50 items was \STwoCEightMed\% lower than tinyBenchmarks' and \STwoCNineMed\% lower than
MINCE's, and over the whole test \STwoCSevenMed\% lower than that of an ATLAS-style 3PL CAT,
in the same direction on every benchmark and significantly on four or five of eight (only the
comparison with MINCE met the frozen criterion). These
methods were designed for other regimes: tinyBenchmarks for hundreds of calibration models,
whose anchors then cover the item space; ATLAS for thousands; MINCE for sizing a random
subset. With a handful of calibrators the plug-in parameters they rely on are unreliable,
and tinyBenchmarks cannot always find as many anchors as its budget. In this regime our evidence
favours adaptive selection under calibration uncertainty, although it is conclusive by the
frozen criterion only against MINCE; with large panels the differences should shrink, as
the gain over plug-in Fisher CAT does.}

\newcommand{\RecommendationsText}{%
For adaptive evaluation of LLMs with a small calibration panel: (1) do not plug in
item-specific discriminations calibrated on fewer than a few dozen systems; integrate their
uncertainty with an empirical-Bayes item prior (the expected response function is cheap),
or constrain them to a common value when a simple, conservative test matters more than
accuracy; (2) select by expected information gain or any information criterion computed
from the marginal model---the choice matters little---and, when the estimand is the
benchmark score, let the score-targeted criterion take over late in the test;
\RecommendationFB\ (4) report model-based intervals with the panel size and their expected
under-coverage, and report coverage alongside accuracy; treat interval recalibration as
promising but not yet validated; (5) exclude the target's whole cluster of related systems, not only its
nominal family, from its panel, and analyse with clusters as units; (6) count the
calibration panel in the cost: with eight fresh calibrators who answer the whole bank,
adaptive testing costs less than full evaluation after \CostBreakFullMin--\CostBreakFullMax\
new systems, but less than uniform subsampling of equal precision only after a median of
\CostBreakNaiveMed; when the calibrators' results already exist, it is cheaper from the first
system.}

\newcommand{\AbstractFB}{Bayesian tests that also integrate the calibrators' abilities were
not significantly different from it,}
\newcommand{\IntroFBFinding}{a posterior-draw marginal Bayesian CAT was not significantly better or worse
than EOST-M on any benchmark but met the frozen equivalence margin on too few of them, before
and after a post hoc correction of its burn-in, and in follow-ups with long chains and
better-mixing samplers the median difference stayed below half a per cent; and}
\newcommand{\DiscussionFBShort}{integrating the panel-level uncertainty as well, in the Bayesian
tests we ran, made no significant difference (equivalence was not shown either).}
\newcommand{\ConclusionFB}{Bayesian CATs that also integrate the calibrators' abilities
were not significantly different from it, without being shown equivalent, and better-mixing
samplers did not change that comparison.}
\newcommand{\RecommendationFB}{(3) the Bayesian treatments of the calibrators that we ran made
no significant difference with panels of this size, so the simpler calibration-marginal model,
exact given the calibrators, is a reasonable default;}
